\documentclass[11pt]{article}
\usepackage[margin=1in]{geometry}
\usepackage{amsmath,amssymb,amsthm,hyperref}
\newtheorem{theorem}{Theorem}
\newcommand{\tr}{\operatorname{tr}}
\title{An exact two-dimensional obstruction to higher-power eigenvalue weak majorization}
\author{Myeongjin Kang}
\date{October 26, 2026}
\begin{document}\maketitle
\begin{abstract}
We verify an explicit counterexample to MathDB problem 372496, namely the universal eigenvalue weak-majorization assertion of Conjecture 5.5 in Zhang's manuscript. All source operators, matrix computations, and the bridge from ordered eigenvalues to trace are certified in Lean 4.34.1. The mathematical counterexample was already posted by Shivam Patel; no discovery or priority claim is made.
\end{abstract}
\noindent\textit{Credit and license.} This AI-assisted paper verifies Shivam Patel's existing mathematical counterexample. Original prose is CC BY 4.0 and original code is MIT.

For Hermitian complex matrices $A,B\in\mathbb H_n$, define exactly as in the source
\[
 H=A+B,\qquad Q_k=\sum_{j=0}^k\binom{k}{j}A^jB^{k-j},\qquad R_k=\frac{Q_k+Q_k^*}{2}.
\]
For decreasing real vectors $x,y\in\mathbb R^n$, weak majorization $x\prec_w y$ means $\sum_{i=1}^r x_i\leq\sum_{i=1}^r y_i$ for every $1\leq r\leq n$. The source asks whether $\lambda(H^k)\prec_w\lambda(R_k)$ for every Hermitian pair and every integer $k\geq5$.
\begin{theorem}
The source assertion fails for $n=2$, $k=5$, and
\[
 A=\begin{pmatrix}0&0\\0&-2\end{pmatrix},\qquad
 B=\begin{pmatrix}0&1\\1&0\end{pmatrix}.
\]
\end{theorem}
\begin{proof}
Both matrices are real symmetric, hence Hermitian as complex matrices. The exact fifth-power and coefficient computations are
\[
 H^5=\begin{pmatrix}-12&29\\29&-70\end{pmatrix},\quad
 Q_5=\begin{pmatrix}0&1\\121&-122\end{pmatrix},\quad
 R_5=\begin{pmatrix}0&61\\61&-122\end{pmatrix}.
\]
For example $H=\left(\begin{smallmatrix}0&1\\1&-2\end{smallmatrix}\right)$ and multiplication yields $H^2=\left(\begin{smallmatrix}1&-2\\-2&5\end{smallmatrix}\right)$; the remaining products are ordinary exact integer arithmetic. The expression for $Q_5$ includes all six binomially weighted terms of its definition, and taking its Hermitian part yields the displayed $R_5$.

For any Hermitian matrix the sum of all eigenvalues, with multiplicities, is its real trace. If the conjectured weak majorization held, its $r=2$ inequality would give
\[
 -82=\tr(H^5)\leq\tr(R_5)=-122,
\]
which is false. This proves failure at an admissible source exponent and hence refutes the universal conjecture.
\end{proof}
The formal module \texttt{WeakMajorization.lean} defines the source $Q_k,R_k$ over complex matrices. It proves every displayed identity and both Hermitian hypotheses. Its predicate \texttt{WeakMajorized} is the initial partial-sum definition (including the harmless empty sum). It applies that predicate to mathlib's actual ordered Hermitian eigenvalues \texttt{eigenvalues₀}; these are decreasing by the spectral theorem. The local trace bridge reindexes mathlib's trace identity to this ordered enumeration. The main theorem is \texttt{WeakMajorizationCounterexample.source\_counterexample}. No spectral data or trace inequality is assumed.

The certificate establishes a full counterexample to Conjecture 5.5. It makes no assertion about the distinct singular-value question $\lambda(H^k)\prec_w\sigma(Q_k)$, and does not formalize the stronger every-odd-order extension included in the posted solution. The certificate is new to this project, not asserted to be new mathematics or the first formalization. Snapshot identifiers and content hashes are supplied in \texttt{source.json}.
\begin{thebibliography}{9}
\bibitem{zhang} T. Zhang, \emph{Weak majorization inequalities for the cubic and quartic coefficients of $e^{(A+B)t}$ versus $e^{At}e^{Bt}$}, arXiv:2601.07286v1, Conjecture 5.5, \url{https://arxiv.org/abs/2601.07286}.
\bibitem{post} S. Patel, posted solution to MathDB 372496, \url{https://mathdb.com/p/372496}. The solution identifier and retrieved content hashes accompany this certificate.
\end{thebibliography}
\end{document}
